\documentclass[12pt]{article}
\providecommand{\palabrasmates}{entero, opuesto, valor absoluto, signo…}
\usepackage{album}
\unidad{Números enteros}
\subtitulo{fuerza y bando · sumar y restar · multiplicar y dividir · potencias}

\begin{document}
\portada

% ═══════════════════════════════ LOS CROMOS ═══════════════════════════════
% Los de la palestra: cada idea, con el campo de la app al lado. Como en el
% álbum de Miut, el andamiaje se retira dentro del guiado: del 1 al 7 casi
% toda la frase está escrita; del 8 en adelante, cada vez menos; el 15, el 16
% y el 17 son solo la pregunta y el espacio.

\debe{\item qué es un número entero: positivos, negativos y el cero \item tres situaciones de la vida con números negativos \item la recta, con el cero en medio}
\begin{cromo}{1}{Los números enteros}{clase · la palestra: el campo y su regla}
\completa Los números enteros son los \h[2.4cm]{positivos} ($+1$, $+2$, $+3$…), los
\h[2.4cm]{negativos} ($-1$, $-2$, $-3$…) y el \h[1.4cm]{cero}.\par\vspace{2mm}
El cero no es ni positivo ni \h[2.4cm]{negativo}.\par\vspace{2mm}
En la palestra, los negativos tiran hacia la \h[2.2cm]{izquierda} y los positivos hacia la \h[2cm]{derecha}.

\vspace{2mm}
\centerline{\campo[paso=4mm]{-3,5}}

\vspace{1mm}
\tuejemplo Tres situaciones de la vida con números negativos:
\espacio[1.4cm]{La temperatura en invierno: $-4$\,°C. El garaje, en la planta $-2$. Una cuenta del banco en números rojos: $-30$\,€.}
\end{cromo}

\debe{\item qué es el valor absoluto, y cómo se escribe \item qué es el signo \item un ejemplo tuyo, con un tirador dibujado}
\begin{cromo}{2}{Fuerza y bando: valor absoluto y signo}{app · Fuerza y bando · Ficha 1}
\begin{minipage}[c]{.72\linewidth}\raggedright
\completa Todo número entero dice dos cosas: \textbf{cuánto} tira y \textbf{hacia dónde}.\par\vspace{2mm}
Cuánto tira es su \h[3.4cm]{valor absoluto}: el número sin su \h[1.4cm]{signo}.
Se escribe entre barras: $|{-5}| =$ \h[1cm]{5}.\par\vspace{2mm}
Hacia dónde tira lo dice su \h[1.6cm]{signo}: $-$ a la izquierda, $+$ a la derecha.\par\vspace{2mm}
En la regla, el valor absoluto es lo lejos que está del \h[1.2cm]{0}.
\end{minipage}%
\begin{minipage}[c]{.28\linewidth}\centering\tirador[1]{-5}\end{minipage}

\vspace{2mm}
\tuejemplo Un número tuyo: su valor absoluto y su signo.
\espacio[1.2cm]{$-7$: su valor absoluto es $|{-7}| = 7$ y su signo es $-$; tira a la izquierda con fuerza 7.}
\end{cromo}

\debe{\item qué son dos números opuestos \item dónde están en la recta \item qué pasa si tiran juntos}
\begin{cromo}{3}{Los opuestos}{app · Fuerza y bando · Ficha 1}
\completa Dos números son \textbf{opuestos} si tienen el mismo \h[3.4cm]{valor absoluto} y distinto \h[1.6cm]{signo}.\par\vspace{2mm}
El opuesto de $-4$ es \h[1.2cm]{$+4$}, y el opuesto de $+9$ es \h[1.2cm]{$-9$}. Se escribe $\text{op}(-4) = +4$.\par\vspace{2mm}
En la regla están a la \h[2.4cm]{misma distancia} del 0, uno a cada lado. Si tiran juntos, la bandera se queda en el \h[1cm]{0}.

\vspace{2mm}
\centerline{\campo[paso=4mm,alcance=6]{-4,4}}

\vspace{1mm}
\tuejemplo Dos opuestos tuyos y su suma:
\espacio[1.1cm]{$-6$ y $+6$: \ $-6 + 6 = 0$.}
\end{cromo}

\debe{\item cómo se sabe, en la recta, cuál es mayor \item por qué $-7$ es menor que $-2$ \item tres números tuyos, ordenados}
\begin{cromo}{4}{Ordenar los enteros}{app · Fuerza y bando (menor y mayor) · Ficha 1}
\completa En la recta, el número que está más a la \h[2.2cm]{izquierda} es el \h[1.6cm]{menor}.\par\vspace{2mm}
Cualquier negativo es \h[1.6cm]{menor} que el 0, y el 0 es menor que cualquier \h[2.2cm]{positivo}.\par\vspace{2mm}
Entre dos negativos, es menor el de \h[3.4cm]{más valor absoluto}: $-7$ \h[.8cm]{$<$} $-2$, aunque el $-7$ tira más fuerte.

\vspace{2mm}
\centerline{\regla[5mm]{8}{-7/izq,-2/izq}}

\vspace{1mm}
\tuejemplo Cuatro números tuyos, ordenados de menor a mayor:
\espacio[1.1cm]{$-6 < -1 < 0 < +3$}
\end{cromo}

\debe{\item qué hace el $+$ en la palestra \item cómo se suman dos enteros del mismo signo \item un ejemplo tuyo con el campo}
\begin{cromo}{5}{Sumar enteros del mismo signo}{app · Entra gente · Ficha 2}
\begin{minipage}[c]{.58\linewidth}\raggedright
\completa El $+$ quiere decir «que \h[1.6cm]{entre}».\par\vspace{2mm}
Si los dos tienen el \textbf{mismo signo}, tiran hacia el mismo lado: se \h[1.6cm]{suman} las fuerzas y se deja el \h[1.4cm]{signo}.\par\vspace{2mm}
$-4 + (-2) =$ \h[1.2cm]{$-6$}
\end{minipage}%
\begin{minipage}[c]{.42\linewidth}\centering\campo[paso=3mm,alcance=7]{-4,-2}\end{minipage}

\vspace{2mm}
\tuejemplo Una suma tuya de dos negativos:
\espacio[1.1cm]{$-3 + (-5) = -8$: tiran los dos a la izquierda, con fuerza $3 + 5 = 8$.}
\end{cromo}

\debe{\item cómo se suman dos enteros de distinto signo \item quién gana \item un ejemplo tuyo con el campo}
\begin{cromo}{6}{Sumar enteros de distinto signo}{app · Entra gente · Ficha 2}
\begin{minipage}[c]{.58\linewidth}\raggedright
\completa Si tienen \textbf{distinto signo}, tiran en contra: se \h[1.6cm]{restan} las fuerzas y gana el signo del que tira \h[2.6cm]{más fuerte}.\par\vspace{2mm}
$-3 + 5 =$ \h[1.2cm]{$+2$}\qquad $4 + (-9) =$ \h[1.2cm]{$-5$}
\end{minipage}%
\begin{minipage}[c]{.42\linewidth}\centering\campo[paso=3mm,alcance=7]{-3,5}\end{minipage}

\vspace{2mm}
\tuejemplo Una suma tuya en la que gane el negativo:
\espacio[1.1cm]{$2 + (-7) = -5$: fuerzas $7 - 2 = 5$, y gana el $-7$, que tira más.}
\end{cromo}

\debe{\item qué hace el $-$ en la palestra \item qué pasa si el que se retira no está \item la regla: restar es sumar el opuesto \item un ejemplo tuyo}
\begin{cromo}{7}{Restar enteros}{app · Alguien se retira · Ficha 2}
\completa El $-$ quiere decir «que se \h[1.8cm]{retire}». Si el que se retira no está, entra con su \h[1.6cm]{pareja} (su opuesto), que juntos suman 0, y se va.\par\vspace{1mm}
{\footnotesize
\begin{tabular}{@{}c@{\hspace{2mm}}c@{\hspace{2mm}}c@{}}
\campo[alcance=6,paso=2.6mm,escala=.45]{2} & \campo[alcance=6,paso=2.6mm,escala=.45]{2,-3,3} & \campo[alcance=6,paso=2.6mm,escala=.45]{2,3}\\
tira el \nm{2} & entran \nm{-3} y \nm{3} & se va el \nm{-3}: \nm{5}\\
\end{tabular}}\par\vspace{2mm}
Por eso \textbf{restar es sumar el} \h[2cm]{opuesto}:\quad $2 - (-3) = 2 +$ \h[1cm]{3} $=$ \h[1cm]{5}

\vspace{2mm}
\tuejemplo Una resta tuya, cambiada a suma:
\espacio[1.1cm]{$-4 - 6 = -4 + (-6) = -10$}
\end{cromo}

\debe{\item cómo se quitan los paréntesis con dos signos seguidos \item los cuatro casos \item un ejemplo tuyo}
\begin{cromo}{8}{Dos signos seguidos}{app · Escríbelo · Ficha 2}
\completa Cuando hay dos signos seguidos, se cambian por uno:\par\vspace{1mm}
{\renewcommand{\arraystretch}{1.6}
\begin{tabular}{@{}l@{\hspace{6mm}}l@{\hspace{6mm}}l@{}}
$+(+3) = +3$ & entra un positivo & tira a la \h[2cm]{derecha}\\
$+(-3) =$ \h[.9cm]{$-3$} & entra un negativo & tira a la \h[2cm]{izquierda}\\
$-(+3) =$ \h[.9cm]{$-3$} & se retira un positivo & la bandera va a la \h[2cm]{izquierda}\\
$-(-3) =$ \h[.9cm]{$+3$} & se retira un negativo & la bandera va a la \h[2cm]{derecha}\\
\end{tabular}}

\vspace{2mm}
\palabras ¿Por qué quitar a uno que tira hacia la izquierda mueve la bandera a la derecha?
\espacio[1.3cm]{Porque deja de tirar hacia la izquierda: lo que queda tira más hacia la derecha que antes.}
\end{cromo}

\debe{\item cómo se hace una cuenta larga de sumas y restas \item un ejemplo tuyo, por bandos}
\begin{cromobrillante}{9}{Sumas y restas encadenadas}{app · Practicar · Ficha 2 reverso}
\completa Primero se quitan los paréntesis: cada resta, a \h[1.6cm]{suma}. Luego se juntan los de cada \h[1.6cm]{bando}:\par\vspace{2mm}
$-5 + 8 - (-3) - 4 = -5 + 8 + 3 - 4$\par\vspace{1mm}
\hspace*{5mm}azules: $8 + 3 =$ \h[1cm]{11}\qquad rojos: $-5 - 4 =$ \h[1cm]{$-9$}\qquad bandera: \h[1cm]{$+2$}

\vspace{2mm}
\tuejemplo Una cuenta tuya de cuatro números, por bandos:
\espacio[1.3cm]{$6 - (-2) + (-7) - 3 = 6 + 2 - 7 - 3$: azules $8$, rojos $-10$: bandera $-2$.}
\end{cromobrillante}

\debe{\item qué hacen los frascos \item la regla de los signos para multiplicar \item un ejemplo tuyo con un frasco dibujado}
\begin{cromo}{10}{Multiplicar enteros: la regla de los signos}{app · Pociones · Ficha 3}
\begin{minipage}[c]{.7\linewidth}\raggedright
\completa Se multiplican las \h[2cm]{fuerzas}, y el signo sale de la regla de los signos:\par\vspace{1mm}
\hspace*{4mm}mismo signo $\rightarrow$ \h[1cm]{$+$}\qquad distinto signo $\rightarrow$ \h[1cm]{$-$}\par\vspace{1mm}
\hspace*{4mm}$(+)\cdot(+) = +$\quad $(-)\cdot(-) =$ \hc{+}\quad $(+)\cdot(-) =$ \hc{−}\quad $(-)\cdot(+) =$ \hc{−}\par\vspace{2mm}
En la palestra, un frasco con calavera ($\cdot(-1)$, $\cdot(-2)$…) le cambia de \h[1.6cm]{bando}: $(-3)\cdot(-2) =$ \h[1cm]{$+6$}
\end{minipage}%
\begin{minipage}[c]{.3\linewidth}\centering\bebe{-3}{\traiciondoble}\end{minipage}

\vspace{2mm}
\tuejemplo Dos productos tuyos, uno que dé positivo y otro negativo:
\espacio[1.1cm]{$(-4)\cdot(-3) = +12$\qquad $5\cdot(-2) = -10$}
\end{cromo}

\debe{\item cómo se dividen dos enteros \item la regla de los signos, también aquí \item un ejemplo tuyo}
\begin{cromo}{11}{Dividir enteros}{app · Pociones (Aguada) · Ficha 3}
\completa Se dividen las \h[2cm]{fuerzas} y el signo sale con la \h[3.6cm]{misma regla} que al multiplicar.\par\vspace{2mm}
$(-12)\dv 3 =$ \h[1cm]{$-4$}\qquad $(-12)\dv(-3) =$ \h[1cm]{$+4$}\qquad $15\dv(-5) =$ \h[1cm]{$-3$}\par\vspace{2mm}
En la palestra, el frasco azul $\dv 2$ deja al tirador en la \h[1.8cm]{mitad} de fuerza, en su bando.

\vspace{2mm}
\tuejemplo Una división tuya con dos negativos:
\espacio[1.1cm]{$(-20)\dv(-4) = +5$}
\end{cromo}

\debe{\item en qué orden se hacen las operaciones \item qué pasa con un menos delante de un paréntesis \item un ejemplo tuyo, paso a paso}
\begin{cromo}{12}{El paréntesis y el orden de las operaciones}{app · El paréntesis · Ficha 3 reverso}
\completa El orden: primero los \h[2.4cm]{paréntesis}, luego las \h[2.4cm]{potencias}, luego multiplicar y \h[1.8cm]{dividir}, y al final sumar y restar.\par\vspace{2mm}
En la palestra, el paréntesis se reduce en el \h[2.2cm]{banquillo} y sale a tirar ya hecho.\par\vspace{2mm}
Un menos delante del paréntesis le cambia de \h[1.6cm]{bando}: $6 - (3 - 8) = 6 - (-5) =$ \h[1.2cm]{11}

\vspace{2mm}
\tuejemplo Una cuenta tuya con paréntesis, paso a paso:
\espacio[1.4cm]{$5 + 2\cdot(-6 + 3) = 5 + 2\cdot(-3) = 5 + (-6) = -1$}
\end{cromo}

\debe{\item qué es una potencia de base negativa \item el signo, según el exponente \item un ejemplo par y uno impar}
\begin{cromo}{13}{Potencias de base negativa}{app · Potencias · Ficha 4}
\begin{minipage}[c]{.72\linewidth}\raggedright
\completa $(-2)^3 = (-2)\cdot(-2)\cdot$\h[1.2cm]{$(-2)$} $=$ \h[1cm]{$-8$}: la base es $-2$ y el exponente, \h[.8cm]{3}.\par\vspace{2mm}
En la palestra, el tirador se bebe \h[3.2cm]{su propio frasco}.\par\vspace{2mm}
Con base negativa, si el exponente es \h[1.4cm]{par} sale positivo; si es \h[1.4cm]{impar}, negativo.
\end{minipage}%
\begin{minipage}[c]{.28\linewidth}\centering\bebe{-2}{\traiciondoble\traiciondoble}\end{minipage}

\vspace{2mm}
\tuejemplo Una potencia tuya con exponente par y otra con impar:
\espacio[1.1cm]{$(-3)^2 = +9$\qquad $(-1)^5 = -1$}
\end{cromo}

\debe{\item qué toca el exponente en $(-3)^2$ y en $-3^2$ \item cuánto vale cada uno}
\begin{cromobrillante}{14}{$(-3)^2$ no es $-3^2$}{app · Potencias · Ficha 4}
\completa El exponente solo afecta a lo que tiene \h[3cm]{justo debajo}.\par\vspace{2mm}
En $(-3)^2$ la base es \h[1cm]{$-3$}, y vale \h[1cm]{$+9$}. En $-3^2$ la base es \h[1cm]{3}, y vale \h[1cm]{$-9$}.

\vspace{2mm}
\palabras ¿Cómo lo explicarías con la palestra?
\espacio[1.4cm]{En $(-3)^2$ el $-3$ se bebe su frasco de $\cdot(-3)$ y vuelve a su bando: $+9$. En $-3^2$ se bebe el $+3$ (sale $+9$) y luego el menos le hace retirarse: $-9$.}
\end{cromobrillante}

\debe{\item el orden de las operaciones, completo \item una cuenta tuya con potencia, paréntesis y frasco}
\begin{cromo}{15}{Todo junto}{app · Todo junto · Ficha 4}
\completa El orden de las operaciones, con un ejemplo:
\espacio[2.2cm]{\paso{1}Paréntesis.\ \paso{2}Potencias.\ \paso{3}Multiplicar y dividir.\ \paso{4}Sumar y restar.\\ $10 - 3\cdot(-2)^2 = 10 - 3\cdot 4 = 10 - 12 = -2$}

\tuejemplo Una cuenta tuya con una potencia, un paréntesis y un producto:
\espacio[1.6cm]{$(-1 + 4)^2 - 2\cdot 5 = 3^2 - 10 = 9 - 10 = -1$}
\end{cromo}

\debe{\item cómo se busca un fallo en una cuenta \item los fallos que más se repiten}
\begin{cromobrillante}{16}{El ojo de Argos}{app · El ojo de Argos · Ficha 4 reverso}
\palabras ¿Cómo encuentras el primer fallo en una cuenta resuelta?
\espacio[1.6cm]{Compruebo línea a línea, desde arriba: cada línea tiene que valer lo mismo que la anterior. La primera que no vale lo mismo es la del fallo.}

\tuejemplo Una cuenta con un fallo, y su corrección:
\espacio[1.6cm]{$-8 - 3\cdot(-2) = -8 - (-6) = -8 - 6$ ✗ \ $\rightarrow$ \ $-8 + 6 = -2$}
\end{cromobrillante}

\debe{\item el error que más has cometido en este tema \item cómo lo evitas}
\begin{cromo}{17}{Mi error típico}{todas las fichas}
\palabras El error que más he cometido en esta unidad, y cómo lo evito:
\espacio[2.6cm]{Este cromo es tuyo: no tiene versión oficial. Un ejemplo de cómo se rellena: «En $-3 - 4$ ponía $+7$, porque “menos por menos es más”. Ahora me pregunto: ¿alguien se retira, o hay un frasco?».}
\end{cromo}

\end{document}
